\documentclass[11pt]{article}
\usepackage[a4paper,margin=1in]{geometry}
\usepackage{amsmath,amsthm,amssymb}
\usepackage{mathtools}

\newtheorem{theorem}{Theorem}[section]
\newtheorem{proposition}[theorem]{Proposition}
\newtheorem{lemma}[theorem]{Lemma}
\newtheorem{corollary}[theorem]{Corollary}
\theoremstyle{definition}
\newtheorem{definition}[theorem]{Definition}
\newtheorem{remark}[theorem]{Remark}

\newcommand{\R}{\mathbb{R}}
\newcommand{\Fr}{\mathcal{A}}
\newcommand{\Om}{\Omega}
\newcommand{\eps}{\varepsilon}
\newcommand{\rstar}{\rho_{*}}
\newcommand{\X}{\mathcal X}
\DeclareMathOperator{\dist}{dist}

\title{Metric Framework and First Variation\\
       of the Volume-Radius Foliation\\
       (Route~$\delta$, Plan~2, Building Block \textsc{MetricFramework})}
\author{}
\date{}

\begin{document}
\maketitle

\begin{abstract}
We set up the metric ambient space $(\X,d_\X)$ of $L^1$-equivalence classes
modulo translations and parametrize the torsion super-level sets by the
volume radius.  The smooth moving-boundary calculation is unconditional and
includes the translation mode in the quotient metric.  The finite-perimeter
metric first-variation estimate is stated in the only form currently justified
by the audit: it is conditional on an upper-gradient recovery theorem giving a
\emph{limsup} approximation of the boundary action.  Standard perimeter lower
semicontinuity and Reshetnyak lower semicontinuity go in the opposite direction
and therefore do not prove the required upper bound for the metric derivative.
The missing recovery theorem must approximate the velocity/action measures of
the level surfaces in the local $L^1$ upper-gradient sense specified below, or
else be replaced by a direct finite-perimeter deformation estimate.  This block
also records that no perimeter-free uniform Lipschitz estimate on bad levels is
proved here; boundedness and nestedness alone are not enough for that claim.
\end{abstract}

\section{Standing hypotheses}\label{sec:setup}

Let $n\ge 2$ and let $\Om\subset\R^n$ be a bounded open set with
$|\Om|=\omega_n=|B_1|$ and $\Om\subset B_R$ (without loss of generality
$R\ge 2$ after rescaling). Let
$u\in H^1_0(\Om)\cap L^\infty(\Om)$ be the torsion potential, the unique
variational solution of
\[
   -\Delta u=1\quad\text{in }\Om,\qquad u=0\quad\text{on }\partial\Om.
\]
The $L^\infty$ bound follows from comparison with the torsion function of
$B_R$; interior elliptic regularity gives $u\in W^{2,p}_{\rm loc}(\Om)$ for
every $p<\infty$. Fix $\rstar\in(0,1)$.
The volume-radius parametrization is
\begin{equation}\label{eq:rho-param}
   E_\rho:=\{u>t(\rho)\},\qquad |E_\rho|=\omega_n\rho^n,\qquad
   \rho\in[\rstar,1],
\end{equation}
with $\rho\mapsto t(\rho)$ absolutely continuous on $[\rstar,1]$ and
$t'(\rho)\le 0$ a.e.\ (by monotonicity of $\rho\mapsto t(\rho)$ and the
volume identity $|E_\rho|=\omega_n\rho^n$; see also \cite{Maggi2012}
Thm.~13.1 and \cite{AFP2000} Thm.~3.40). For a.e.\ $\rho$, $E_\rho$ has finite
perimeter, reduced boundary $\partial^*E_\rho$, outer unit normal $\nu_\rho$,
and $P(E_\rho)=\mathcal H^{n-1}(\partial^*E_\rho)$. We write
$P_\rho:=P(B_\rho)=n\omega_n\rho^{n-1}$.

\section{The metric space $\X$ and the rescaled curve}\label{sec:X}

\begin{definition}[$L^1$ modulo translations]\label{def:X}
Let $\mathfrak M$ denote the set of $\mathcal L^n$-equivalence classes of
characteristic functions $\mathbf 1_A$ with $|A|<\infty$, equipped with the
$L^1$ distance $d_{L^1}(\mathbf 1_A,\mathbf 1_C)=|A\Delta C|$. The group
$\R^n$ acts on $\mathfrak M$ by translation. The metric space
$(\X,d_\X)$ is the quotient
\begin{equation}\label{eq:X-def}
   \X:=\mathfrak M/\R^n,\qquad
   d_\X([A],[C]):=\inf_{a\in\R^n}|A\Delta(C+a)|.
\end{equation}
\end{definition}

\begin{lemma}[Attainment of the translation infimum]\label{lem:attain}
If $A,C\subset\R^n$ have finite measure and are both contained in some ball
$B_M$, the infimum in~\eqref{eq:X-def} is attained at some
$a_*\in\overline{B_{2M}}$.
\end{lemma}

\begin{proof}
For $|a|\ge 2M$, $(C+a)\cap A=\emptyset$ and $|A\Delta(C+a)|=|A|+|C|$, which
exceeds $|A|+|C|-|A\cap C|=|A\Delta C|+|A\cap C|\ge|A\Delta C|$, so the
infimum is unchanged if we restrict to $a\in\overline{B_{2M}}$. The map
$a\mapsto|A\Delta(C+a)|$ is continuous on $\R^n$ (by dominated convergence),
hence attains its minimum on the compact set $\overline{B_{2M}}$.
\end{proof}

\begin{definition}[Rescaled curve]\label{def:F}
For $\rho\in[\rstar,1]$, define the volume-rescaled level set and its class
\[
   F_\rho:=\rho^{-1}E_\rho\subset\R^n,\qquad [F_\rho]\in\X.
\]
\end{definition}

By construction $|F_\rho|=\omega_n=|B_1|$ for all $\rho$. Since
$E_\rho\subset B_R$, we have $F_\rho\subset B_{R/\rstar}$ uniformly in
$\rho\in[\rstar,1]$, and Lemma~\ref{lem:attain} applies to all pairs
$([F_{\rho_1}],[F_{\rho_2}])$.

\begin{definition}[Metric derivative; Ambrosio--Gigli--Savar\'e
{\cite[Thm.~1.1.2]{AGS2008}}]\label{def:AC}
A curve $\gamma:I\to\X$ is absolutely continuous if there is $m\in L^1(I)$
with $d_\X(\gamma(s),\gamma(\sigma))\le\int_\sigma^s m$ for
$\sigma\le s$ in $I$. The metric derivative
\[
   |\dot\gamma|(\rho):=\lim_{h\to 0}\frac{d_\X(\gamma(\rho+h),\gamma(\rho))}{|h|}
\]
exists for a.e.\ $\rho\in I$ and is the smallest admissible $m$.
\end{definition}

\section{The normal velocity $V_\rho$ and the Sobolev--Sard step}\label{sec:V}

For a.e.\ $\rho\in[\rstar,1]$, set
\begin{equation}\label{eq:V-def}
   V_\rho(x):=\frac{-t'(\rho)}{|\nabla u(x)|}\qquad
   \text{for }x\in\partial^*E_\rho,\ |\nabla u|(x)>0.
\end{equation}

\begin{proposition}[Sobolev--Sard for $W^{2,p}_{\rm loc}$]\label{prop:sard}
Let $p>n$ (via Morrey embedding $W^{2,p}\hookrightarrow C^{1,\alpha}$,
$\alpha=1-n/p$). Then
\begin{equation}\label{eq:sard}
   \mathcal H^{n-1}\bigl(\partial^*E_\rho\cap\{|\nabla u|=0\}\bigr)=0
   \qquad\text{for a.e.\ }\rho\in[\rstar,1].
\end{equation}
\end{proposition}

\begin{proof}
By Stampacchia, $u\in W^{2,p}_{\rm loc}(\Om)$ for every $p<\infty$; fix any
$p>n$. The Sobolev--Sard theorem of Figalli and Maggi
\cite[Appendix~A, Thm.~A.1]{FigalliMaggi2011} states that for $u\in
W^{2,p}_{\rm loc}$ with $p>n$, one has
$\mathcal H^{n-1}(\{u=t\}\cap\{|\nabla u|=0\})=0$ for a.e.\
$t\in(0,\|u\|_\infty)$. By the De Giorgi structure theorem,
$\partial^*\{u>t\}\subset\{u=t\}$ up to an $\mathcal H^{n-1}$-null set for
a.e.\ $t$ (Maggi \cite[\S 15]{Maggi2012}; \cite{AFP2000} Thm.~3.59).
Reparametrizing by $\rho$ via~\eqref{eq:rho-param} yields~\eqref{eq:sard}.
\end{proof}

\begin{remark}[Backup citations]\label{rem:backup}
The same statement is recorded in De Philippis--Marini--Mukoseeva 2021,
Lemma~2.4, with a self-contained proof. The classical Brothers--Ziemer
1988, Lemma~2.4, gives a $W^{2,p}$ Sard statement modulo a
capacity-to-$\mathcal H^{n-1}$ conversion.
\end{remark}

\begin{corollary}[$V_\rho$ is pointwise defined and $L^1$]\label{cor:V-L1}
For a.e.\ $\rho\in[\rstar,1]$, $V_\rho$ is an $\mathcal H^{n-1}$-a.e.\
finite, strictly positive Borel function on $\partial^*E_\rho$, and
\begin{equation}\label{eq:V-flux}
   \int_{\partial^*E_\rho}V_\rho\,d\mathcal H^{n-1}
   =\frac{d}{d\rho}|E_\rho|=n\omega_n\rho^{n-1}=:A_\rho.
\end{equation}
In particular $V_\rho\in L^1(\partial^*E_\rho,\mathcal H^{n-1})$.
\end{corollary}

\section{Centre fields}\label{sec:center}

\begin{definition}[Borel centre field]\label{def:center}
A \emph{centre field} for $\{E_\rho\}$ is a Borel-measurable map
$z:[\rstar,1]\to\R^n$ with $\sup_\rho|z_\rho|\le R'$ for some
$R'<\infty$ that is any radius containing all translates considered.
For each $\rho\in[\rstar,1]$ and $x\in\partial^*E_\rho$,
the associated \emph{homothetic velocity} is
\begin{equation}\label{eq:H-def}
   H_{z_\rho,\rho}(x):=\frac{(x-z_\rho)\cdot\nu_\rho(x)}{\rho}.
\end{equation}
\end{definition}

\begin{remark}[Downstream centre choices]
\label{rem:centroid}
This building block treats $z_\rho$ as an abstract Borel centre field.
The repaired metric trace block instantiates $z_\rho$ as a Fusco--Julin
oscillation centre on good levels.  The centroid field of
\cite{CentroidBound} remains useful for the auxiliary space--time $\Pi$
identity, but it is not needed for the metric closure.
\end{remark}

\section{The metric first-variation theorem (T)}\label{sec:T}

\begin{theorem}[Conditional metric first variation modulo translations]\label{thm:T}
Under the standing hypotheses of Section~\ref{sec:setup} and for any Borel
centre field $z=\{z_\rho\}$ with $\sup_\rho|z_\rho|\le R'<\infty$, assume the
upper-gradient recovery theorem in Definition~\ref{def:UGR}.  Then the curve
$\rho\mapsto[F_\rho]$ is absolutely continuous from $[\rstar,1]$ into
$(\X,d_\X)$, and
\begin{equation}\label{eq:T}
   |\dot F_\rho|_\X
   \le \rho^{-n}\inf_{a\in\R^n}
       \int_{\partial^*E_\rho}\bigl|V_\rho-H_{z_\rho,\rho}-a\cdot\nu_\rho\bigr|
       \,d\mathcal H^{n-1}
\end{equation}
for a.e.\ $\rho\in[\rstar,1]$.
\end{theorem}

\subsection{Smooth flows: the moving-region identity}\label{ss:T-smooth}

Suppose temporarily that $\Om$ is smooth, $u\in C^\infty(\overline{\Om})$,
and $|\nabla u|>0$ on the strip
$\bigcup_{\rho\in[\rstar,1-\eta]}\partial E_\rho$ for some $\eta>0$.  In the
pointwise calculation below $z\in\R^n$ is frozen at the base value
$z=z_\rho$; no derivative of the centre field is used.  The homothety
\[
   \Psi_s^z(y):=z+s^{-1}(y-z)
\]
maps $E_s$ onto $z+s^{-1}(E_s-z)$, which represents the same point of $\X$ as
$F_s=s^{-1}E_s$.

\begin{lemma}[Pulled-back normal velocity]\label{lem:W}
In the smooth case, fix $z=z_\rho$ and a translation velocity
$b\in\R^n$ in the rescaled variables.  The normal velocity of
\[
   s\longmapsto \Psi_s^z(E_s)-(s-\rho)b
\]
at $s=\rho$, pulled back to $\partial E_\rho$ through $\Psi_\rho^z$, equals
\begin{equation}\label{eq:W}
   \rho^{-1}\bigl(V_\rho-H_{z,\rho}\bigr)-b\cdot\nu_\rho
   =
   \rho^{-1}\bigl(V_\rho-H_{z,\rho}-a\cdot\nu_\rho\bigr),
   \qquad a:=\rho b,
   \qquad\text{on }\partial E_\rho.
\end{equation}
\end{lemma}

\begin{proposition}[Smooth case of (T)]\label{prop:T-smooth}
In the smooth setting above,
\begin{equation}\label{eq:T-smooth-X}
   \limsup_{h\to 0}\frac{d_\X([F_{\rho+h}],[F_\rho])}{|h|}
   \le \rho^{-n}\inf_{a\in\R^n}\int_{\partial^*E_\rho}
       \bigl|V_\rho-H_{z_\rho,\rho}-a\cdot\nu_\rho\bigr|\,d\mathcal H^{n-1}.
\end{equation}
\end{proposition}

The proof is the first-variation identity \cite[\S17.3]{Maggi2012}
applied to the translated rescaled flow in Lemma~\ref{lem:W}.  The boundary
Jacobian of $\Psi_\rho^z$ contributes $\rho^{1-n}$ and the normal velocity
has the extra factor $\rho^{-1}$, giving $\rho^{-n}$.  Taking the infimum over
the rescaled translation velocity $b$, equivalently over $a=\rho b$, gives
\eqref{eq:T-smooth-X}.

\subsection{Finite-perimeter passage: required upper-gradient recovery}
\label{ss:T-bv}

The previous draft tried to pass from Proposition~\ref{prop:T-smooth} to
finite-perimeter level sets by perimeter lower semicontinuity, Reshetnyak
lower semicontinuity, and lower semicontinuity of metric derivatives.  That
argument proves the wrong inequality for this purpose.  If $G_k$ denotes the
smooth boundary action and $G$ the target action, Reshetnyak gives
$G\le\liminf_k G_k$, while Theorem~\ref{thm:T} needs an upper gradient bounded
by $G$.  A valid finite-perimeter closure must therefore be a genuine
\emph{limsup} recovery theorem, or a direct deformation theorem for
finite-perimeter sets.

\begin{definition}[Upper-gradient recovery needed for Theorem~T]\label{def:UGR}
Fix a bounded Borel centre field $z=\{z_\rho\}$ and set
\begin{equation}\label{eq:G-target}
   G_z(\rho):=\rho^{-n}\inf_{a\in\R^n}
       \int_{\partial^*E_\rho}
       |V_\rho-H_{z_\rho,\rho}-a\cdot\nu_\rho|\,d\mathcal H^{n-1}.
\end{equation}
The \emph{upper-gradient recovery theorem} is the following statement.
There exist smooth nested families $\{E_\rho^k\}_{\rho\in[\rstar,1]}$ with
volume-radius parametrisations and curves
$\gamma_k(\rho):=[(E_\rho^k)/\rho]\in\X$ such that:
\begin{enumerate}
\item $\gamma_k(\rho)\to\gamma(\rho):=[F_\rho]$ in $\X$ for every
      $\rho\in[\rstar,1]$, after changing the representative of the original
      level curve on a null set if necessary;
\item the smooth boundary actions
\[
   G_{z,k}(\rho):=\rho^{-n}\inf_{a\in\R^n}
       \int_{\partial^*E_\rho^k}
       |V_\rho^k-H^k_{z_\rho,\rho}-a\cdot\nu_\rho^k|\,
       d\mathcal H^{n-1}
\]
      are upper gradients for $\gamma_k$;
\item for every interval $I\subset[\rstar,1]$,
\begin{equation}\label{eq:UGR}
   \limsup_{k\to\infty}\int_I G_{z,k}(\rho)\,d\rho
   \le
   \int_I G_z(\rho)\,d\rho .
\end{equation}
\end{enumerate}
\end{definition}

\begin{remark}[What must be proved elsewhere]
Condition~\eqref{eq:UGR} is the missing velocity-measure convergence in the
correct direction.  It is not a consequence of strict $BV$ convergence plus
Reshetnyak lower semicontinuity.  A proof would need a recovery construction
for the full boundary action, including the scalar normal velocity
$V_\rho=-t'(\rho)/|\nabla u|$, the homothetic term
$H_{z_\rho,\rho}$, the normals, and the translation infimum.  Equivalently,
one may bypass recovery by proving directly that for a.e.\ $\rho$ and every
$a\in\R^n$ the finite-perimeter deformation generated by the normal speed
$V_\rho-H_{z_\rho,\rho}-a\cdot\nu_\rho$ gives the one-sided quotient-metric
upper bound in \eqref{eq:T}.
\end{remark}

\begin{proof}[Proof of Theorem~\ref{thm:T}]
Assume Definition~\ref{def:UGR}.  For $s<t$, pointwise convergence in $\X$
and the smooth upper-gradient estimate give
\[
   d_\X(\gamma(s),\gamma(t))
   =\lim_{k\to\infty}d_\X(\gamma_k(s),\gamma_k(t))
   \le \limsup_{k\to\infty}\int_s^t G_{z,k}(\rho)\,d\rho
   \le \int_s^t G_z(\rho)\,d\rho .
\]
Thus $G_z\in L^1_{\rm loc}$ is an upper gradient for $\gamma$, so $\gamma$ is
absolutely continuous and the metric derivative satisfies
$|\dot\gamma|(\rho)\le G_z(\rho)$ for a.e.\ $\rho$ by
Definition~\ref{def:AC}.  This is exactly~\eqref{eq:T}.
\end{proof}

\section{Translations quotient kills the $a\cdot\nu_\rho$ mode}\label{sec:trans}

\begin{lemma}[Translation derivative formula]\label{lem:trans-deriv}
For every bounded set $E\subset\R^n$ of finite perimeter and every
$a\in\R^n$,
\begin{equation}\label{eq:trans-deriv}
   \lim_{\eps\to 0}\frac{|E\Delta(E+\eps a)|}{\eps}
   =\int_{\partial^*E}|a\cdot\nu_E|\,d\mathcal H^{n-1}.
\end{equation}
(\cite{Maggi2012} Thm.~13.8; \cite{AFP2000} Thm.~3.84.)
\end{lemma}

\section{Per-$\rho$ squared form: status and open issues}\label{sec:open}

\subsection{Pointwise (per-$\rho$) version: conditional}

If one assumes the per-$\rho$ radial trace bound (statement~(R), to be
supplied by Building Block \cite{WeightedMetricTrace}; audit-trail markdown
note) and the homothetic velocity defect bound,
\begin{equation}\label{eq:25}
   \Bigl(\int_{\partial^*E_\rho}|H_{z_\rho,\rho}-\bar V_\rho|
   \,d\mathcal H^{n-1}\Bigr)^2\le C_{n,\rstar,R}D_I(t(\rho)),
\end{equation}
then Theorem~\ref{thm:T}, the triangle inequality, and Cauchy--Schwarz
combine to give the pointwise per-$\rho$ squared bound
\begin{equation}\label{eq:71}
   |\dot F_\rho|_\X^2\le C_{n,\rstar,R}\bigl(D_H(t(\rho))+D_I(t(\rho))\bigr).
\end{equation}

\begin{remark}[Honest status]\label{rem:status}
The estimate \eqref{eq:25} is conditional on the pointwise (R), which is
open in the no-graph setting. Consequently \eqref{eq:71} is \emph{conditional}
and not used downstream.
\end{remark}

\subsection{Integrated version: repaired route}

The Plan~2 paper aims to close the bound in a $\mu$-integrated good-level
form.  With Definition~\ref{def:UGR} or a direct finite-perimeter deformation
theorem in place, Theorem~\ref{thm:T} can be combined on the
small-isoperimetric-defect levels with the levelwise Fusco--Julin centre, the
oriented radial trace estimate, and the self-truncated velocity estimate.
This metric block does \emph{not} provide a perimeter-free uniform Lipschitz
bound on the complementary bad levels.  Such a bound cannot be inferred merely
from boundedness of $E_\rho\subset B_R$ and nestedness of the level sets,
because the dilation part may carry uncontrolled perimeter.  Bad levels must
therefore be handled either by a separate proved estimate from the level-set
structure or by an argument that avoids pointwise bad-level Lipschitz control.

\section{Summary and use downstream}\label{sec:summary}

\begin{theorem}[Statement of \textsc{MetricFramework}]\label{thm:main}
Let $n\ge 2$, $0<\rstar<1$, and $R\ge 2$. Let $\Om\subset B_R$ be bounded
open with $|\Om|=\omega_n$, let $u$ be its torsion potential, let $\{E_\rho\}$
be the volume-radius foliation~\eqref{eq:rho-param}, and let
$F_\rho:=\rho^{-1}E_\rho$ in $\X$. For every Borel centre field $z=\{z_\rho\}$
with $\sup_\rho|z_\rho|\le R'$, the following statements hold.
\begin{itemize}
\item the normal velocity $V_\rho=-t'(\rho)/|\nabla u|$ is $\mathcal H^{n-1}$-a.e.\
finite and positive on $\partial^*E_\rho$ for a.e.\ $\rho$;
\item in the smooth noncritical setting, the quotient first-variation bound
holds with the sharp scaling factor $\rho^{-n}$ as in
Proposition~\ref{prop:T-smooth};
\item in finite-perimeter generality, if the upper-gradient recovery theorem
of Definition~\ref{def:UGR} is proved for this centre field, then
$\rho\mapsto[F_\rho]\in\X$ is absolutely continuous and
\[
   |\dot F_\rho|_\X\le \rho^{-n}\inf_{a\in\R^n}
   \int_{\partial^*E_\rho}\bigl|V_\rho-H_{z_\rho,\rho}-a\cdot\nu_\rho\bigr|
   \,d\mathcal H^{n-1}\qquad\text{for a.e.\ }\rho\in[\rstar,1].
\]
\end{itemize}
\end{theorem}

\begin{thebibliography}{99}

\bibitem{AGS2008}
L.~Ambrosio, N.~Gigli, G.~Savar\'e,
\emph{Gradient Flows in Metric Spaces and in the Space of Probability
Measures}, 2nd ed., Birkh\"auser, 2008.

\bibitem{AFP2000}
L.~Ambrosio, N.~Fusco, D.~Pallara,
\emph{Functions of Bounded Variation and Free Discontinuity Problems},
Oxford, 2000.

\bibitem{Maggi2012}
F.~Maggi, \emph{Sets of Finite Perimeter and Geometric Variational
Problems}, Cambridge, 2012.

\bibitem{FigalliMaggi2011}
A.~Figalli, F.~Maggi, ARMA \textbf{201} (2011), 143--207. App.~A, Thm.~A.1.

\bibitem{LevelSetIdentity}
\emph{Building Block} \textsc{LevelSetIdentity}, Plan~2.

\bibitem{CentroidBound}
\emph{Building Block} \textsc{CentroidBound}, Plan~2.

\bibitem{SpaceTimeIdentity}
\emph{Building Block} \textsc{SpaceTimeIdentity}, Plan~2.

\bibitem{WeightedMetricTrace}
\emph{Building Block} \textsc{WeightedMetricTrace}, Plan~2.

\end{thebibliography}

\end{document}
